\documentclass[10pt,a4paper]{amsart}
\usepackage[margin=19mm]{geometry}
\usepackage{amsmath,amssymb,mathtools}
\usepackage[scr=rsfs,cal=euler]{mathalfa}
\usepackage{xfrac}
\usepackage[unicode,hidelinks,bookmarks=false]{hyperref}
\newcommand{\ie}{i.e.\ }
\newcommand{\cf}{cf.\ }
\newcommand{\resp}{respectively\ }
\newcommand{\Subsection}{\subsection}
\providecommand{\nobreakdash}{\nobreak\hbox{-}}
\setlength{\emergencystretch}{2em}
\multlinegap0pt
\usepackage{enumitem}
\newcommand{\R}{\mathbb{R}}
\newcommand{\B}{\mathscr{B}}
\newcommand{\ind}{\mathbf{1}}
\newcommand{\n}{\hspace*{-5pt}}
\DeclareMathOperator{\poisson}{Poisson}
\DeclareMathOperator{\uniform}{Uniform}
\DeclareMathOperator{\binomial}{Binomial}
\DeclareMathOperator{\card}{card}
\newtheorem{theorem}{Theorem}
\renewenvironment{proof}{\noindent{\bf Proof.}}{\hspace*{2mm}~$\square$}
\title[Conditioning property for multi-dimensional Poisson point processes]{Short proof of the conditioning property for multi-dimensional Poisson point processes: the conditioning property (Theorem 2)}
\author{Nicolas Lanchier}
\date{Source: arXiv:2508.21303v1 (CC BY 4.0)}
\begin{document}
\begin{abstract}
Poisson processes and one-dimensional Poisson point processes satisfy three main properties: superposition, thinning, and conditioning.
The proof of the first two relies on basic estimates involving the Poisson distribution that are also true for multi-dimensional Poisson point processes.
In contrast, the proof of conditioning uses that the distances between consecutive occurrences in time or entities in space are independent and exponentially distributed, which is nonsensical in higher dimensions.
This paper gives a short proof of the conditioning property for multi-dimensional Poisson point processes.
\end{abstract}
\maketitle

\noindent\emph{Source and licence.} Short proof of the conditioning property for multi-dimensional Poisson point processes. Version arXiv:2508.21303v1 (CC BY 4.0), distributed under the Creative Commons Attribution 4.0 International licence, http://creativecommons.org/licenses/by/4.0/, as recorded in the arXiv OAI licence metadata for this exact version and shown on the abstract page https://arxiv.org/abs/2508.21303v1. Full rights notice: licenses/arxiv-2508.21303-CC-BY-4.0.txt.\par\smallskip
\noindent\textit{Editorial scope.} Counted unit: the paper's conditioning property for multi-dimensional Poisson point processes, printed as Theorem 2 (source0.tex lines 113--122, unlabelled in the source, carrying the optional title Conditioning property), delivered together with the paper's complete original proof (source0.tex lines 123--157) as the single primary proof body. No mathematics is written, repaired or re-derived: statement and proof are the paper's own text line for line, in the paper's own notation ($\R$, $\B$, $\ind$, $\poisson$, $\uniform$, $\binomial$, $\card$, $\lambda$, $\phi$). The paper renews the proof environment as a run-in heading with a square, and that renewed environment at source0.tex lines 123--157 is exactly the primary proof span; the wrapper reproduces the paper's renewal verbatim. Required context retained uncounted: the whole of Section 1 with the definition of the Poisson process through axioms (A1) and (A2) and of the $d$-dimensional Poisson point process through axioms (A1') and (A2'), the identification of the one-dimensional Poisson point process with the Poisson process, the description of superposition, thinning and conditioning, and the statement of the one-dimensional conditioning property (printed Theorem 1, source0.tex lines 90--98) that the counted theorem extends; and the whole of Section 2 down to the definition of the uniform random variable on a Borel set given by its density (source0.tex lines 59--112). Every definition, theorem, numbered display, footnote and citation of those lines is retained, including the two citations that record where the classical one-dimensional proof may be found and the citation of the author's video explanation that introduces the counted theorem. Omitted from this package and not redistributed: the paper's front matter (source0.tex lines 32--55: its titlepage block, address, AMS subject classification, keywords and dedication, of which the title and the author name are reproduced in the wrapper title block and the abstract is reproduced verbatim in the wrapper front matter), the paper's own bibliography listing beyond the cited entries, and the paper's closing rule lines. Nothing else is omitted: apart from its front matter, its references and the two sections retained here the paper has no further body content, so the delivered document contains the whole of the paper's mathematics. The delivered text cites all three of the paper's bibliography keys (lanchier\_2017, ross\_2010 and lanchier\_2024, four citation commands in all), and those three entries are transcribed verbatim from the paper's own in-source thebibliography (source0.tex lines 161--173) in the paper's own order, so the printed bibliography numbers of the delivered unit are the paper's own (1, 2 and 3). Reference adaptation: none was needed and none was made. Every reference inside the delivered unit resolves inside it: the label th:poisson of the retained one-dimensional Theorem 1, which Section 2 cites in order to say that the counted theorem extends it, and the labels eq:uniform, eq:poisson1 and eq:poisson2 of the three numbered displays inside the counted proof, which the proof's own two closing cross-references cite. No unresolved reference or citation is produced, and the two theorems keep the paper's own numbering, by which the one-dimensional statement is Theorem 1 and the counted statement is Theorem 2. Printed slips kept literal: no printed word, symbol, hypothesis or number of the counted statement or of its proof was altered. The paper's own forms are delivered as printed, among them its use of the script B for the Borel sigma-algebra, its notation N(B) for the count of Poisson points in a Borel set, its use of card for cardinality, its habit of writing probabilities with a capital P and of conditioning with a vertical bar, the exact wording of axiom (A1') with the word disjoint, its own spacing macros inside the proof displays, its optional-argument theorem title Conditioning property, and its own phrase in distribution. The three numbered displays inside the counted proof are numbered (1), (2) and (3) in the delivered unit exactly as in the published version, because they are the paper's first three numbered equations as well. Non-ASCII characters: the paper contains no character outside ASCII, so no encoding substitution, no glyph replacement and no missing-character risk arises. Layout and class: the paper's own class is article with the options letterpaper, 11pt, twoside, keywordsasfootnote, addressatend and noinfoline, and its frontmatter, its theorem numbering and its page layout are supplied by the locally shipped imsart.sty that travels in the e-print archive; imsart.sty is NOT redistributed with this package and is not used to build it. The wrapper is the lane's pinned amsart editorial wrapper at 10pt on A4, to which this package adds the paper's own macro block (its definitions for R, the script B, the indicator, the negative-space macro and the four operators Poisson, Uniform, Binomial and card), the paper's own theorem declaration and the paper's own renewal of the proof environment, all copied verbatim, plus the enumitem package that the retained itemize lists with explicit left margins require. The paper's own loads of imsart, fullpage and babel are therefore replaced by the wrapper's pinned geometry and fonts; that is a page-layout difference of the wrapper only, and no formula, symbol, hypothesis or argument changes. No publisher class, no publisher style file and no upstream script is redistributed. Assets: the paper has no figure environment, no graphics-inclusion command, no PDF-page inclusion and no table; no retained span contains a graphics-inclusion command, an input directive or an include directive; the package declares no asset dependency, and no image file is copied into or redistributed with this package. All mathematics of the unit is therefore native text with no image dependency. A rights notice is delivered at licenses/arxiv-2508.21303-CC-BY-4.0.txt.
\section{Poisson and Poisson point processes}
 The Poisson random variable can be viewed as the limit of binomial random variables as the number of trials~$n$ goes to infinity and the success probability scales like~$1/n$, making this random variable the continuous-time analog of the binomial random variable.
 Poisson processes and Poisson point processes are derived from the Poisson random variable, and are used respectively to model the distribution of occurrences in time and the distribution of entities in space.
 In particular, Poisson point processes are extensions of Poisson processes to potentially higher dimensions.
 More precisely, the process~$X = \{X_t : t \in [0, \infty) \}$ is the Poisson process with rate~$\mu$ if it starts at~$X_0 = 0$ and if it has independent Poisson increments, meaning that
\begin{itemize}[leftmargin=36pt]
\item[(A1)] For all~$t_1 < t_2 < \cdots < t_{2k - 1} < t_{2k}$, $X_{t_2} - X_{t_1}, \ldots, X_{t_{2k}} - X_{t_{2k - 1}}$ are independent. \vspace*{4pt}
\item[(A2)] For all~$s < t$, $X_t - X_s = \poisson (\mu (t - s))$.
\end{itemize}
 The random variable~$X_t$ is interpreted as the number of occurrences until time~$t$.
 This concept can be extended to higher dimensions using a counting process
 $$ N = \{N (B) : B \in \B (\R^d) \ \hbox{bounded} \} \quad \hbox{where} \quad \B (\R^d) = \hbox{Borel~$\sigma$-algebra} $$
 and where~$N (B)$ is interpreted as the number of entities in the Borel set~$B$.
 The process~$N$ is the Poisson point process with intensity~$\mu$ if the number of entities, called Poisson points, in nonoverlapping subsets are independent, and the number of entities in a given subset is Poisson distributed.
 More precisely, letting~$\lambda$ refer to the Lebesgue measure,
\begin{itemize}[leftmargin=40pt]
\item[(A1')] For all~$B_1, B_2, \ldots, B_k \in \B (\R^d)$ disjoint, $N (B_1), N (B_2), \ldots, N (B_k)$ are independent. \vspace*{4pt}
\item[(A2')] For all~$B \in \B (\R^d)$ bounded, $N (B) = \poisson (\mu \lambda (B))$.
\end{itemize}
 Note that the Poisson point process~$N$ on the real line~(for~$d = 1$) and the Poisson process~$X$ are mathematically equivalent in the sense that
 $$ X_t = N ([0, t]) \quad \hbox{in distribution \ for all \ $t \geq 0$}. $$
 These two processes satisfy three main properties, called superposition, thinning, and conditioning, that can be described for the one-dimensional Poisson point process~$N$ as follows. \vspace*{4pt} \\
\noindent{\bf Superposition}.
 Having a palette of~$n$ colors, labeled~$i = 1, 2, \ldots, n$, assume that entities with color~$i$ are randomly distributed on the real line according to independent Poisson point processes with intensity~$\mu_i$.
 Then, the set of entities of any color~(what you see if you become colorblind) is distributed according to a Poisson point process with intensity~$\mu_1 + \mu_2 + \cdots + \mu_n$. \vspace*{4pt} \\
\noindent{\bf Thinning}.
 Having a palette of~$n$ colors, labeled~$i = 1, 2, \ldots, n$, and entities randomly distributed on the real line according to a Poisson point process with intensity~$\mu$, paint each entity independently color~$i$ with probability~$p_i$ where~$p_1 + p_2 + \cdots + p_n = 1$.
 Then, the~$n$ sets of entities with different colors are distributed according to independent Poisson point processes with intensity~$\mu p_i$. \vspace*{4pt} \\
\noindent{\bf Conditioning}.
 Having entities randomly distributed on the real line according to a Poisson point process with intensity~$\mu$, given that there are~$n$ entities in~$[0, t]$, the position of these entities is described by~$n$ independent uniform random variables on the interval~$[0, t]$. \vspace*{4pt} \\
 From the point of view of the Poisson process~$X$, this can be formulated as
\begin{theorem}[Conditioning property]
\label{th:poisson}
 Given that~$X_t = n$,
\begin{itemize}
\item let $U_1, U_2, \ldots, U_n$ be independent~$\uniform (0, t)$, \vspace*{4pt}
\item let $\tau_i = \inf \,\{s : X_s = i \}$ for~$i = 1, 2, \ldots, n$ be the times of the occurrences in~$(0, t)$.
\end{itemize}
 Then~$\{U_1, U_2, \ldots, U_n \} = \{\tau_1, \tau_2, \ldots, \tau_n \}$ in distribution.
\end{theorem}
\noindent
 Superposition and thinning can be viewed as converse of each other.
 Their proof relies on basic estimates involving the Poisson random variable that are not sensitive to the dimension, so they are also true for multi-dimensional Poisson point processes.
 In contrast, the main ingredient to prove the conditioning property is that the distances between consecutive occurrences in time or entities in space are independent and exponentially distributed with parameter~$\mu$, a property that is nonsensical in higher dimensions~(see~\cite[Lemma~9.4]{lanchier_2017} or~\cite[Theorem~5.2]{ross_2010} for a proof).
 Also, though the conditioning property is expected to hold for multi-dimensional Poisson point processes, the traditional proof in one dimension does not extend to higher dimensions, and the author could not find any proof in the probability literature.
 The main objective of this paper is to give a short proof of the conditioning property for multi-dimensional Poisson point processes.
 
%%%%%%%%%%%%%%%%%%%%%%%%%%%%%%%%%%%%%%%%%%%%%%%%%%%%%%%%%%%%%%%%%%%%%%%%%%%%%%%%%%%%%%%%%%%%%%%%%%%%%%%%%%%%%%%%%%%%%%%%%%%%%%%%%%%%%%%%%%%%%%%%%%%%%%%%%%%%%%%%%%%%%%%%%%%%%%%%%%%%%%%%%%%%%%%%%%

\section{The Conditioning property}
 To extend Theorem~\ref{th:poisson} to higher dimensions, we let~$N$ be the~$d$-dimensional Poisson point process with intensity~$\mu$, and we let~$B$ be a bounded Borel subset of~$\R^d$.
 The uniform random variable on this Borel set is the continuous random variable with density function
 $$ \phi : \R^d \to \R_+ \quad \hbox{defined as} \quad \phi = \ind_B / \lambda (B). $$
 Then, we have the following result that the author also explains in the video~\cite{lanchier_2024}.

\begin{theorem}[Conditioning property]
 Given that~$N (B) = n$,
\begin{itemize}
\item let $U_1, U_2, \ldots, U_n$ be independent~$\uniform (B)$, \vspace*{4pt}
\item let $P_1, P_2, \ldots, P_n \in \R^d$ be the Poisson points in~$B$.
\end{itemize}
 Then~$\{U_1, U_2, \ldots, U_n \} = \{P_1, P_2, \ldots, P_n \}$ in distribution in the sense that
 $$ P (\card (A \cap \{U_1, U_2, \ldots, U_n \}) = k) = P (\card (A \cap \{P_1, P_2, \ldots, P_n \}) = k) $$
 for all Borel subsets~$A \subset B$ and all~$k = 0, 1, \ldots, n$.
\end{theorem}

\begin{proof}
 Let~$a = \lambda (A)$ and~$b = \lambda (B)$.
 By definition of the uniform random variable,
 $$ \begin{array}{c} P (U_i \in A) = \int \phi \,\ind_A \,d\lambda = \lambda (A) / \lambda (B) = a/b \quad \hbox{for all} \quad i = 1, 2, \ldots, n. \end{array} $$
 In particular, using also independence, for all~$k = 0, 1, \ldots, n$,
\begin{equation}
\label{eq:uniform}
\begin{array}{rcl}
  P (\card (A \cap \{U_1, U_2, \ldots, U_n \}) = k) & \n = \n & P (\card \{i = 1, 2, \ldots, n : U_i \in A \} = k) \vspace*{4pt} \\
                                                    & \n = \n & P (\binomial (n, a/b) = k). \end{array}
\end{equation}
 Now, because~$A \cap (B \setminus A) = \varnothing$, it follows from axiom~(A1') that
\begin{equation}
\label{eq:poisson1}
\begin{array}{l}
  P (\card (A \cap \{P_1, P_2, \ldots, P_n \}) = k \,| \,N (B) = n) \vspace*{4pt} \\ \hspace*{20pt} =
  P (\card \{i = 1, 2, \ldots, n : P_i \in A \} = k \,| \,N (B) = n) \vspace*{4pt} \\ \hspace*{20pt} =
  P (N (A) = k \,| \,N (B) = n) \vspace*{4pt} \\ \hspace*{20pt} =
  P (N (A) = k, N (B) = n) / P (N (B) = n) \vspace*{4pt} \\ \hspace*{20pt} =
  P (N (A) = k, N (B \setminus A) = n - k) / P (N (B) = n) \vspace*{4pt} \\ \hspace*{20pt} =
  P (N (A) = k) P (N (B \setminus A) = n - k) / P (N (B) = n). \end{array}
\end{equation}
 Using also that~$\lambda (B \setminus A) = \lambda (B) - \lambda (A) = b - a$ and axiom~(A2') of the definition of a Poisson point process, then simplifying, the last term in~\eqref{eq:poisson1} becomes
\begin{equation}
\label{eq:poisson2}
\begin{array}{l}
% \displaystyle \frac{P (\poisson (\mu a) = k) \,P (\poisson (\mu (b - a)) = n - k)}{P (\poisson (\mu b) = n)} \vspace*{12pt} \\ \hspace*{40pt} =
\displaystyle \bigg(\frac{\mu^k a^k \,e^{- \mu a}}{k!} \bigg) \bigg(\frac{\mu^{n - k} (b - a)^{n - k} \,e^{- \mu (b - a)}}{(n - k)!} \bigg) \bigg/ \bigg(\frac{\mu^n b^n \,e^{- \mu b}}{n!} \bigg) \vspace*{12pt} \\ \hspace*{40pt} =
\displaystyle \bigg(\frac{a^k}{k!} \bigg) \bigg(\frac{(b - a)^{n - k}}{(n - k)!} \bigg) \bigg/ \bigg(\frac{b^n}{n!} \bigg) =
\displaystyle \frac{n!}{(n - k)! \,k!} \ \bigg(\frac{a^k (b - a)^{n - k}}{b^k \,b^{n - k}} \bigg) \vspace*{12pt} \\ \hspace*{40pt} =
\displaystyle {n \choose k} \bigg(\frac{a}{b} \bigg)^k \bigg(1 - \frac{a}{b} \bigg)^{n - k} =
\displaystyle P (\binomial (n, a/b) = k). \end{array}
\end{equation}
 The result follows noticing that the last terms in~\eqref{eq:uniform} and~\eqref{eq:poisson2} are equal.
\end{proof}

\begin{thebibliography}{99}
\bibitem[1]{lanchier_2017} N. Lanchier. \emph{Stochastic Modeling}. Universitext. Springer, Cham, 2017. xiii + 303pp.
\bibitem[2]{lanchier_2024} N. Lanchier (Oct 12, 2024).
Chapter 09. Poisson processes and the exponential distribution~(with subtitles) [Video]. YouTube. \href{https://www.youtube.com/watch?v=1A3wlvXcBjE}{https://www.youtube.com/watch?v=1A3wlvXcBjE}
\bibitem[3]{ross_2010} S. M. Ross. \emph{Introduction to Probability Models}, 10th ed. (Academic Press, 2010).
\end{thebibliography}
\end{document}
